% Propietario conservado: /Users/ruben/Documents/New project/output/REVISION_ARGUMENTAL_APERTURAS_CIERRES_20260915/III/spanish_source/manuscrito/sections/11_nucleo_finito.tex
\section{Enumeración exhaustiva y determinación finita de los operadores de transporte}
\label{iii:iii:nucleo-finito}

Este apéndice desarrolla dos operaciones finitas con datos de partida
diferentes. La primera enumera las condiciones iniciales de los dos cursores
y construye sus emisiones mediante la recurrencia aritmética declarada.
La segunda determina los operadores de transporte a partir del registro
regional de transiciones. La separación de sus dominios permite precisar
qué produce cada demostración: la enumeración construye el catálogo de
emisiones; la reconstrucción matricial determina los operadores compatibles
con el registro. El teorema de reconstrucción no se emplea como demostración
de la producción anterior de ese registro.

Las primitivas de la primera operación son el soporte $I_9^2$, las cuatro
orientaciones cardinales, el selector trítico, la lectura anterior al
desplazamiento, las inversiones de orientación y la regla de emisión.
Sus definiciones se reúnen a continuación. Las de la segunda son las
cuatro matrices de transiciones de~\eqref{iii:iii:nf-registro}.
Su reconstrucción única se demuestra en el
teorema~\ref{iii:iii:nf-unicidad}; el calendario resultante se expresa
en~\eqref{iii:iii:nf-calendario}. Estas determinaciones conservan
los datos de partida declarados para cada operación.

\subsection{Dominio completo de condiciones iniciales}

Sea $I_9=\{0,\ldots,8\}$ y sea
\[
 \mathcal D=\{N,E,S,O\},\qquad
 v_N=(-1,0),\quad v_E=(0,1),\quad
 v_S=(1,0),\quad v_O=(0,-1).
\]
Una condición inicial de cursor es $s=(i,j,d)\in
\mathcal S=I_9^2\times\mathcal D$. Los dos cursores son independientes
en el dominio inicial
\begin{equation}
 \Omega=\mathcal S_+\times\mathcal S_\times,
 \qquad |\mathcal S|=9^2\cdot4=324,
 \qquad |\Omega|=324^2=104\,976.
 \label{iii:iii:nf-dominio}
\end{equation}
El índice $i$ corresponde a la etiqueta positiva $i+1$ de APP. Para
enteros positivos, la reducción positiva es $\rho_9(n)=1+[(n-1)]_9$,
donde $[a]_b$ denota el representante en $\{0,\ldots,b-1\}$.
Las evaluaciones residuales de las dos hojas son
\[
 \sigma(i,j)=\rho_9(i+j+2),\qquad
 \varpi(i,j)=\rho_9((i+1)(j+1)).
\]
El lector multiplicativo finito utiliza
\begin{equation}
 \begin{array}{c|rrrrrrrrr}
 a&1&2&3&4&5&6&7&8&9\\ \hline
 \ell_9(a)&0&1&0&2&5&0&4&3&0
 \end{array}.
 \label{iii:iii:nf-log}
\end{equation}
Sobre las unidades, $\ell_9$ es el exponente discreto en base $2$.
Sobre $3,6,9$, el valor $0$ constituye la extensión declarada del
observable de emisión; la pertenencia a ese sector permanece en el
registro de trayectoria.

Para $1\le t\le54$, defínase
\[
 r_t=1+[(t-1)]_9,
 \qquad
 \tau_t=
 \begin{cases}
 +1,&r_t\in\{1,4,7\},\\
 -1,&r_t\in\{2,5,8\},\\
 0,&r_t\in\{3,6,9\}.
 \end{cases}
\]
Cada cursor conserva su posición $z_t$ inmediatamente antes del paso
$t$ y su vector de orientación $v_t$. El cursor aditivo está activo
cuando $\tau_t=+1$; el multiplicativo, cuando $\tau_t=-1$. Para cada
uno de ellos, con indicador de actividad $a_t$, la actualización es
\begin{equation}
 z_{t+1}=[z_t+a_tv_t]_9,
 \qquad
 v_{t+1}=
 \begin{cases}
 -v_t,&t\in\{27,54\},\\
 v_t,&t\notin\{27,54\}.
 \end{cases}
 \label{iii:iii:nf-actualizacion}
\end{equation}
La lectura del paso $t$ se realiza en $z_t$, antes del desplazamiento.
El evento neutral conserva ambas posiciones. Esta convención fija
inequívocamente los extremos de las ventanas y el efecto de las
inversiones cardinales.

En las ventanas $W_m=\{9m-8,\ldots,9m\}$, $1\le m\le6$, sean
\begin{align}
 S_m(s_+)&=\sum_{\substack{t\in W_m\\\tau_t=+1}}
                 \sigma(z_t^+),&
 E_m(s_\times)&=\sum_{\substack{t\in W_m\\\tau_t=-1}}
                 \ell_9(\varpi(z_t^\times)),
 \label{iii:iii:nf-firmas}\\
 s_m&=[S_m]_{10},&e_m&=[E_m]_{10}.\notag
\end{align}
En cada ventana hay tres eventos de cada tipo. Las firmas $s$ y $e$
son las seis componentes de estos residuos, no los enteros $S_m,E_m$
anteriores a su reducción.

\subsection{Censo de firmas y de emisiones}

La emisión de una pareja de condiciones iniciales queda determinada por
\begin{equation}
 U_m=100s_m+10[s_m+e_m]_{10}+[3s_m+5e_m+1]_{10},
 \qquad 1\le m\le6.
 \label{iii:iii:nf-emision}
\end{equation}
El término final $1$ es $[7\cdot3]_{10}$, correspondiente a los tres
eventos neutrales. Los coeficientes de esta regla forman parte de la
recurrencia declarada. Las constantes reales que se estudian posteriormente
no intervienen en la enumeración de su dominio ni en su evaluación.

\begin{proposition}[Factorización inyectiva del catálogo finito]
\label{iii:iii:nf-inyectividad}
La emisión distingue cada pareja de firmas $(s,e)$. Hay $18$ firmas
aditivas y $26$ multiplicativas; sus pares producen $468$ emisiones.
Las multiplicidades de estas emisiones son $144$, $432$ y $1944$,
con $432$, $18$ y $18$ emisiones respectivamente.
\end{proposition}

\begin{proof}
La cifra de centenas de $U_m$ recupera $s_m$. Si $b_m$ es su cifra de
decenas, entonces $e_m=[b_m-s_m]_{10}$. La cifra de unidades verifica
la tercera congruencia de \eqref{iii:iii:nf-emision}. Por tanto, el
acoplamiento es inyectivo.

El censo se obtiene recorriendo, en el orden cartesiano completo,
$i=0,\ldots,8$, $j=0,\ldots,8$ y $d=N,E,S,O$, y aplicando los $54$
pasos de \eqref{iii:iii:nf-actualizacion}. Las $324$ condiciones aditivas
se distribuyen en las siguientes $18$ firmas, cada una con $18$
preimágenes. En las tablas, una sucesión de seis cifras representa
las seis componentes ordenadas de una firma.
\begin{center}
\begin{tabular}{ccc}
$122564$&$122898$&$212645$\\
$212988$&$221456$&$221889$\\
$456221$&$465898$&$546988$\\
$564122$&$645212$&$654889$\\
$889221$&$889654$&$898122$\\
$898465$&$988212$&$988546$
\end{tabular}
\end{center}
Para el cursor multiplicativo, las firmas y sus multiplicidades son
\begin{center}
\begin{tabular}{rr@{\qquad}rr}
Firma&Multiplicidad&Firma&Multiplicidad\\ \hline
$000000$&$108$&$555555$&$24$\\
$177393$&$8$&$555717$&$8$\\
$177555$&$8$&$555771$&$8$\\
$177771$&$8$&$555933$&$8$\\
$339555$&$8$&$717555$&$8$\\
$339771$&$8$&$717717$&$8$\\
$339933$&$8$&$717933$&$8$\\
$393177$&$8$&$771177$&$8$\\
$393393$&$8$&$771339$&$8$\\
$393555$&$8$&$771555$&$8$\\
$555177$&$8$&$933339$&$8$\\
$555339$&$8$&$933555$&$8$\\
$555393$&$8$&$933717$&$8$
\end{tabular}
\end{center}
Estas tablas resultan de las sumas finitas \eqref{iii:iii:nf-firmas}, con
dominio, orden de lectura e inversiones ya especificados. Sus sumas
de multiplicidades son $18\cdot18=324$ y
$24\cdot8+24+108=324$. El dominio cartesiano
\eqref{iii:iii:nf-dominio} permite realizar cualquier pareja de firmas.
La inyectividad proporciona $18\cdot26=468$ emisiones y las
multiplicidades
\[
 (18\cdot24,\ 18\cdot8)=(432,144),\qquad
 (18,\ 18\cdot24)=(18,432),\qquad
 (18,\ 18\cdot108)=(18,1944),
\]
donde cada par indica número de emisiones y multiplicidad de su fibra.
Finalmente,
\[
 432\cdot144+18\cdot432+18\cdot1944=104\,976.
\]
La enumeración agota así el dominio completo de condiciones iniciales.
\end{proof}

La reducción ternaria orientada es
\[
 \Psi(U)=([-U_1]_3,\ldots,[-U_6]_3).
\]
Aplicada a las $18\cdot26$ emisiones de las tablas anteriores,
produce el siguiente censo de fibras. Aquí $r$ cuenta emisiones
distintas, no condiciones iniciales.
\begin{equation}
 \begin{array}{c|rrrrrr}
 r&1&2&3&4&5&6\\ \hline
 \#\{w:\#\Psi^{-1}(w)=r\}&132&66&18&3&6&18
 \end{array}.
 \label{iii:iii:nf-censo-ternario}
\end{equation}
La suma de la segunda fila es $243$, mientras que su suma ponderada
por $r$ es $468$. Se distingue, por ello, la imagen de $243$ palabras
del ambiente $\mathbb F_3^6$, que contiene $729$ palabras. El certificado
adjunto conserva todas las firmas con sus preimágenes, las $468$
emisiones y su reducción ternaria; las fórmulas anteriores proporcionan
un procedimiento completo para reconstruirlas.

\subsection{Registro de transiciones y reconstrucción de los levantamientos}

En este apartado se fija el registro regional del propietario citado.
Sean $b_0^\chi,\ldots,b_4^\chi\in\mathbb F_3^6$, con
$\chi\in\{\mathsf C,\mathsf P,\mathsf A\}$, sus tres sucesiones.
El registro organiza las transiciones $b_0\to b_1\to b_2$ en
$X_0,Y_0$ y las transiciones $b_2\to b_3\to b_4$ en $X_1,Y_1$,
con vectores fila y el mismo orden regional:
\begingroup\small
\begin{equation}
\begin{aligned}
X_0&=\begin{pmatrix}
0&1&0&2&1&1\\0&1&2&2&2&2\\2&0&1&1&0&1\\
1&2&1&2&2&1\\1&2&1&2&0&0\\1&1&2&2&0&2
\end{pmatrix},&
Y_0&=\begin{pmatrix}
0&1&2&2&2&2\\0&1&0&2&1&1\\1&2&1&2&2&1\\
1&0&2&0&1&1\\1&1&2&2&0&2\\1&2&1&0&2&0
\end{pmatrix},\\[1ex]
X_1&=\begin{pmatrix}
0&1&0&2&1&1\\0&0&2&1&1&1\\1&0&2&0&1&1\\
0&1&2&2&2&2\\1&2&1&0&2&0\\0&1&0&2&1&0
\end{pmatrix},&
Y_1&=\begin{pmatrix}
0&0&2&1&1&1\\1&1&0&2&2&1\\0&1&2&2&2&2\\
1&0&2&0&1&1\\0&1&0&2&1&0\\0&1&0&2&0&0
\end{pmatrix}.
\end{aligned}
\label{iii:iii:nf-registro}
\end{equation}
\endgroup
La producción precedente se expresa en su propietario mediante la
composición de selección, transporte y actualización:
\[
 x_{j+1}^{\chi}
 =\mathcal U_{9j+9}\circ\cdots\circ\mathcal U_{9j+1}(x_j^{\chi}),
 \qquad b_j^\chi=\Pi_6(x_j^\chi),
 \qquad\mathcal U_t=\mathsf{Upd}_t\circ\mathsf{Tra}_t\circ\mathsf{Sel}_t.
\]
El resultado siguiente parte del registro \eqref{iii:iii:nf-registro}.
Su demostración permite reconstruir los operadores y verificar la
compatibilidad de las transiciones; no atribuye a la inversión matricial
la producción coordenada de los estados $x_j^\chi$.

\begin{proposition}[Unicidad sobre el registro regional]
\label{iii:iii:nf-unicidad}
Los sistemas $X_aL_a=Y_a$, $a=0,1$, tienen una única solución sobre
$\mathbb F_3$. Las dos soluciones son
\begingroup\small
\begin{equation}
L_0=\begin{pmatrix}
2&2&2&1&2&1\\2&1&2&2&1&1\\1&0&0&1&0&2\\
0&0&1&1&1&0\\2&2&2&0&2&1\\2&1&2&1&0&0
\end{pmatrix},\qquad
L_1=\begin{pmatrix}
2&1&2&0&2&2\\1&2&1&1&2&0\\1&2&1&1&1&2\\
0&1&0&0&2&1\\2&0&2&1&0&1\\0&2&2&2&1&1
\end{pmatrix}.
\label{iii:iii:nf-lifts}
\end{equation}
\endgroup
\end{proposition}

\begin{proof}
La eliminación en $\mathbb F_3$ da $\det X_0=1$ y
$\det X_1=2=-1$. Cada aplicación $L\mapsto X_aL$ es, por tanto,
un automorfismo del espacio de matrices de dimensión $36$; en una
vectorización por columnas está representada por $I_6\otimes X_a$.
Se obtiene $L_a=X_a^{-1}Y_a$. La multiplicación de las matrices de
\eqref{iii:iii:nf-lifts} por las respectivas matrices $X_a$ reproduce las
dos matrices $Y_a$ de \eqref{iii:iii:nf-registro}; existencia y unicidad
quedan establecidas sobre ese registro.
\end{proof}

Las tres sucesiones resultantes, con separación explícita en bloques,
son
\begin{align*}
 B^{\mathsf C}&=010211\mid012222\mid010211\mid002111\mid110221,\\
 B^{\mathsf P}&=201101\mid121221\mid102011\mid012222\mid102011,\\
 B^{\mathsf A}&=121200\mid112202\mid121020\mid010210\mid010200.
\end{align*}
Cada bloque conserva sus seis posiciones. La concatenación no se
identifica con un entero en base diez.

\subsection{Exhaustividad del calendario y estabilidad de la reconstrucción}

Un calendario binario es $c=(c_0,c_1,c_2,c_3)\in\{0,1\}^4$.
Defínase el número de transiciones incompatibles con el registro por
\begin{equation}
 d(c)=\sum_{\chi\in\{\mathsf C,\mathsf P,\mathsf A\}}
       \sum_{j=0}^{3}
       \mathbf1\{b_j^\chi L_{c_j}\ne b_{j+1}^\chi\}.
 \label{iii:iii:nf-calendario}
\end{equation}
Las multiplicaciones de las sucesiones anteriores por
\eqref{iii:iii:nf-lifts} dan
\[
 d(c)=3\,d_{\rm H}(c,(0,0,1,1)),
\]
donde $d_{\rm H}$ es la distancia de Hamming: en cada una de las
cuatro posiciones, el operador alternativo incumple las tres
transiciones regionales. El censo completo es
\begin{center}
\begin{tabular}{cr@{\qquad}cr@{\qquad}cr@{\qquad}cr}
$c$&$d(c)$&$c$&$d(c)$&$c$&$d(c)$&$c$&$d(c)$\\ \hline
$0000$&$6$&$0100$&$9$&$1000$&$9$&$1100$&$12$\\
$0001$&$3$&$0101$&$6$&$1001$&$6$&$1101$&$9$\\
$0010$&$3$&$0110$&$6$&$1010$&$6$&$1110$&$9$\\
$0011$&$0$&$0111$&$3$&$1011$&$3$&$1111$&$6$
\end{tabular}
\end{center}
En consecuencia, $0011$ es el único calendario que reproduce las
tres sucesiones completas.

La sensibilidad a las entradas de las matrices también admite una
prueba exhaustiva sin aproximaciones. Una modificación unitaria
sustituye $L_a$ por $L_a+\lambda E_{ij}$, con
$a\in\{0,1\}$, $1\le i,j\le6$ y
$\lambda\in\mathbb F_3\setminus\{0\}$. Hay
$2\cdot6\cdot6\cdot2=144$ modificaciones de este tipo. Como
$X_a$ es invertible,
\[
 X_a(L_a+\lambda E_{ij})-Y_a
 =\lambda X_aE_{ij}\ne0.
\]
Toda modificación destruye al menos una de las transiciones del
registro correspondiente. El certificado ejecutable enumera además
las $144$ modificaciones y conserva, para cada una, el número de
transiciones incumplidas.

La propiedad demostrada es la rigidez de los operadores respecto del
registro regional fijado. Para reproducir una construcción anterior
del propio registro se necesita la acción coordenada de sus operadores
de selección, transporte y actualización. Esa dependencia se conserva
separada del teorema de rigidez, del catálogo finito de emisiones y
de las realizaciones analíticas posteriores.

\subsection{Certificado de reproducción y residencia de los datos}

El programa \texttt{verificar\_nucleo\_finito.py} recorre las $324$
condiciones de cada cursor y forma el producto de las clases resultantes.
El catálogo documental se carga después de esta construcción para
comparar las emisiones y sus multiplicidades. En la segunda parte,
las cuatro matrices de transiciones se extraen de la copia literal
del propietario; las soluciones se calculan por eliminación exacta
módulo $3$ antes de compararlas con las matrices impresas. El recibo
\texttt{NUCLEO\_FINITO.json} conserva los dominios, las preimágenes,
las emisiones, los $16$ calendarios y las $144$ modificaciones.

La comprobación utiliza enteros y residuos exactos. No recibe valores
metrológicos ni expansiones reales de las constantes. Su alcance
es el de las operaciones finitas y los datos de partida especificados
en este apéndice. Las demostraciones de compatibilidad a profundidad
ilimitada y las identificaciones físicas conservan sus respectivos
dominios y obligaciones expositivas.
